\documentclass{article}
\usepackage{amsmath,amssymb,amsthm}
\usepackage{algorithm}
\usepackage{algorithmic}
\usepackage{bm}
\usepackage{enumitem}
\usepackage{hyperref}
\newtheorem{theorem}{Theorem}
\newtheorem{lemma}{Lemma}
\newtheorem{assumption}{Assumption}
\newcommand{\E}{\mathbb{E}}
\newcommand{\R}{\mathbb{R}}
\newcommand{\norm}[1]{\left\|#1\right\|}
\newcommand{\ip}[2]{\langle #1,#2\rangle}
\begin{document}

\section*{Randomized Coordinate Descent with Probability Weighting (RCD-PW)}

\section*{Mathematical Formulation}
Let $f:\R^{d}\rightarrow\R$ be a differentiable convex function.  
For each coordinate $i\in\{1,\dots,d\}$ we are given a selection probability
\[
p_i>0,\qquad \sum_{i=1}^{d}p_i=1 .
\]
Define the \emph{coordinate gradient estimator}
\[
\boxed{\; g^{(t)} \;:=\; \frac{1}{p_{i_t}}\; e_{i_t}\,\partial_{i_t} f(w^{(t)})\;}
\tag{1}
\]
where $i_t$ is drawn independently from $\{1,\dots,d\}$ according to $p$, and
$e_i$ denotes the $i$‑th unit vector.  
The update rule of the algorithm is
\[
\boxed{\; w^{(t+1)} \;=\; w^{(t)} \;-\; \eta\, g^{(t)}\;}
\tag{2}
\]
with a fixed stepsize $\eta>0$ (or a stepsize chosen from a diminishing schedule).  
Because of the factor $1/p_{i_t}$, the estimator $g^{(t)}$ satisfies
\[
\E\bigl[g^{(t)}\mid w^{(t)}\bigr]=\nabla f\!\bigl(w^{(t)}\bigr).
\tag{3}
\]

\section*{Pseudocode}
\begin{algorithm}[H]
\caption{Randomized Coordinate Descent with Probability Weighting}
\begin{algorithmic}[1]
\STATE \textbf{Input:} initial point $w^{(0)}\in\R^{d}$, stepsize $\eta>0$, probabilities $\{p_i\}_{i=1}^{d}$.
\FOR{$t=0,1,2,\dots,T-1$}
    \STATE Sample $i_t\sim\text{Multinomial}(p_1,\dots,p_d)$
    \STATE Compute partial derivative $g_{i_t}^{(t)}:=\partial_{i_t}f\!\bigl(w^{(t)}\bigr)$
    \STATE Form unbiased estimator $g^{(t)}:=\frac{1}{p_{i_t}}e_{i_t} g^{(t)}_{i_t}$
    \STATE Update $w^{(t+1)}:= w^{(t)} - \eta\, g^{(t)}$
\ENDFOR
\STATE \textbf{Output:} $w^{(T)}$ (or the average $\bar w_T:=\frac1T\sum_{t=0}^{T-1}w^{(t)}$)
\end{algorithmic}
\end{algorithm}

\section*{Dry Run Example}
Consider the one‑dimensional quadratic
\[
f(w)=(w-3)^2,
\qquad \nabla f(w)=2(w-3).
\]
Here $d=1$, $p_1=1$, $\eta=0.1$, $w^{(0)}=0$.

\begin{tabular}{c|c|c|c}
$t$ & $i_t$ & $w^{(t+1)}=w^{(t)}-\eta\frac{1}{p_{i_t}}\partial_{i_t}f$ & $f\!\bigl(w^{(t+1)}\bigr)$\\\hline
0 & 1 & $0-0.1\cdot 2(0-3)=0+0.6=0.6$ & $(0.6-3)^2=5.76$\\
1 & 1 & $0.6-0.1\cdot 2(0.6-3)=0.6-0.1\cdot(-4.8)=1.08$ & $(1.08-3)^2=3.6864$\\
2 & 1 & $1.08-0.1\cdot 2(1.08-3)=1.08-0.1\cdot(-3.84)=1.464$ & $(1.464-3)^2=2.3689$
\end{tabular}

The objective value decreases at each iteration, illustrating the algorithm’s progress.

\newpage
\section*{Assumptions}
\begin{assumption}[Convexity]\label{as:convex}
$f$ is convex, i.e.
\[
f(y)\ge f(x)+\ip{\nabla f(x)}{y-x},\qquad \forall x,y\in\R^{d}.
\]
\end{assumption}
\begin{assumption}[Coordinate-wise $L_i$‑smoothness]\label{as:smooth}
For each $i$ there exists $L_i>0$ such that for all $x\in\R^{d}$ and any scalar $\alpha$,
\[
\bigl|\partial_i f(x+\alpha e_i)-\partial_i f(x)\bigr|\le L_i|\alpha|.
\]
Equivalently,
\[
f(x+\alpha e_i)\le f(x)+\alpha\partial_i f(x)+\frac{L_i}{2}\alpha^{2}.
\tag{4}
\]
\end{assumption}
\begin{assumption}[Strong convexity (optional)]\label{as:strong}
$f$ is $\mu$‑strongly convex ($\mu>0$):
\[
f(y)\ge f(x)+\ip{\nabla f(x)}{y-x}+\frac{\mu}{2}\|y-x\|^{2},\qquad\forall x,y.
\]
\end{assumption}
\begin{assumption}[Stepsize]\label{as:step}
The stepsize satisfies $0<\eta\le \frac{1}{\max_i L_i}$ (for the non‑strongly‑convex case) or
$0<\eta\le \frac{2}{\mu+\max_i L_i}$ (for the strongly convex case).
\end{assumption}

\section*{Main Convergence Theorem}
\begin{theorem}[Expected suboptimality]\label{thm:conv}
Let Assumptions \ref{as:convex}–\ref{as:step} hold.
Define $L_{\max}:=\max_i L_i$ and the weighted norm
\[
\|x\|_{P^{-1}}^{2}:=\sum_{i=1}^{d}\frac{x_i^{2}}{p_i},
\qquad 
\|g\|_{P}^{2}:=\sum_{i=1}^{d}p_i g_i^{2}.
\]
Then for the iterates generated by (2) we have for any $T\ge1$:
\[
\boxed{\;
\E\!\bigl[f(w^{(T)})-f(w^{\star})\bigr]
\;\le\;
\frac{\|w^{(0)}-w^{\star}\|_{P^{-1}}^{2}}{2\eta T}
\;}
\tag{5}
\]
where $w^{\star}$ is a minimizer of $f$.
If additionally $f$ is $\mu$‑strongly convex (Assumption \ref{as:strong}), then
\[
\boxed{\;
\E\!\bigl[f(w^{(T)})-f(w^{\star})\bigr]
\;\le\;
\bigl(1-\eta\mu\bigr)^{T}\,\bigl[f(w^{(0)})-f(w^{\star})\bigr]
\;}
\tag{6}
\]
provided $\eta\le 1/L_{\max}$.
\end{theorem}

\section*{Theoretical Properties}
\begin{itemize}[leftmargin=*]
\item \textbf{Unbiasedness via weighting.} The factor $1/p_{i_t}$ in (1) guarantees
$\E[g^{(t)}\mid w^{(t)}]=\nabla f(w^{(t)})$, a crucial property for the proof.
\item \textbf{Stability w.r.t.\ probabilities.} Larger $p_i$ reduces the variance
contribution of the $i$‑th coordinate because the term $\tfrac{1}{p_i}$ appears only
inside the unbiased estimator, not in the expectation.
\item \textbf{Comparison to uniform RCD.} Setting $p_i=1/d$ recovers the classic
randomized coordinate descent bound $\mathcal{O}(d/L_{\max}T)$. By choosing
$p_i\propto L_i^{-1/2}$ one can improve the constant factor.
\item \textbf{Hyperparameter sensitivity.} The stepsize bound \eqref{as:step}
depends only on the largest coordinate smoothness $L_{\max}$, not on the
probability vector; thus the algorithm is robust to the choice of $\{p_i\}$ as
long as they remain strictly positive.
\end{itemize}

\section*{Discussion}
The theorem shows that, despite updating only a single coordinate per iteration,
the algorithm attains the same $\mathcal{O}(1/T)$ convergence rate (in expectation) as full‑gradient methods for convex smooth objectives. In the strongly convex regime the linear rate (6) mirrors that of gradient descent, with the contraction factor $1-\eta\mu$ independent of $d$. The weighting by $1/p_i$ is essential: without it the estimator would be biased and the proof would break down.

\newpage
\section*{Preliminaries (Definitions and Lemmas)}
\begin{lemma}[Coordinate smoothness inequality]\label{lem:smooth}
For any $w\in\R^{d}$, any coordinate $i$, and any scalar $\alpha$,
\[
f(w+\alpha e_i)\le f(w)+\alpha\,\partial_i f(w)+\frac{L_i}{2}\alpha^{2}.
\]
\end{lemma}
\begin{proof}
Directly from Assumption \ref{as:smooth}. \hfill $\square$
\end{proof}

\begin{lemma}[Unbiased gradient estimator]\label{lem:unbiased}
Let $g^{(t)}$ be defined in \eqref{1}. Then
\[
\E\!\bigl[g^{(t)}\mid w^{(t)}\bigr]=\nabla f\!\bigl(w^{(t)}\bigr).
\]
\end{lemma}
\begin{proof}
\[
\E\!\bigl[g^{(t)}\mid w^{(t)}\bigr]
=\sum_{i=1}^{d}p_i\frac{1}{p_i}e_i\partial_i f(w^{(t)})
=\sum_{i=1}^{d}e_i\partial_i f(w^{(t)})=\nabla f(w^{(t)}).
\] \hfill $\square$
\end{proof}

\section*{Main Proof (Step‑by‑Step Derivation)}
We prove Theorem \ref{thm:conv}. Throughout the proof we condition on the sigma‑field
$\mathcal{F}_t:=\sigma\!\bigl(w^{(0)},i_0,\dots,i_{t-1}\bigr)$.

\bigskip
\noindent\textbf{[Step 1]}  Using the update rule (2) and expanding the weighted norm
$\|w^{(t+1)}-w^{\star}\|_{P^{-1}}^{2}$:
\[
\begin{aligned}
\|w^{(t+1)}-w^{\star}\|_{P^{-1}}^{2}
&= \Bigl\|w^{(t)}-w^{\star}-\eta g^{(t)}\Bigr\|_{P^{-1}}^{2}\\
&= \|w^{(t)}-w^{\star}\|_{P^{-1}}^{2}
   -2\eta\ip{g^{(t)}}{w^{(t)}-w^{\star}}_{P^{-1}}
   +\eta^{2}\|g^{(t)}\|_{P^{-1}}^{2}.
\end{aligned}
\tag{7}
\]

\noindent\textbf{[Step 2]}  Compute the inner product term. By definition of the
$P^{-1}$‑weighted inner product,
\[
\ip{g^{(t)}}{w^{(t)}-w^{\star}}_{P^{-1}}
= \sum_{i=1}^{d}\frac{1}{p_i}g^{(t)}_i\bigl(w^{(t)}_i-w^{\star}_i\bigr).
\]
Since $g^{(t)}$ has only the sampled coordinate $i_t$ non‑zero,
\[
\ip{g^{(t)}}{w^{(t)}-w^{\star}}_{P^{-1}}
= \frac{1}{p_{i_t}}\partial_{i_t}f(w^{(t)})\bigl(w^{(t)}_{i_t}-w^{\star}_{i_t}\bigr).
\tag{8}
\]

\noindent\textbf{[Step 3]}  Compute the norm term $\|g^{(t)}\|_{P^{-1}}^{2}$.
Again only coordinate $i_t$ survives:
\[
\|g^{(t)}\|_{P^{-1}}^{2}
= \sum_{i=1}^{d}\frac{1}{p_i}\bigl(g^{(t)}_i\bigr)^{2}
= \frac{1}{p_{i_t}}\Bigl(\frac{1}{p_{i_t}}\partial_{i_t}f(w^{(t)})\Bigr)^{2}
= \frac{1}{p_{i_t}^{2}}\partial_{i_t}f(w^{(t)})^{2}.
\tag{9}
\]

\noindent\textbf{[Step 4]}  Take conditional expectation of \eqref{7} w.r.t.\ $\mathcal{F}_t$.
Using Lemma \ref{lem:unbiased} and the facts that $p_{i_t}$ is the sampling
probability,
\[
\begin{aligned}
\E\!\bigl[\|w^{(t+1)}-w^{\star}\|_{P^{-1}}^{2}\mid\mathcal{F}_t\bigr]
=&\; \|w^{(t)}-w^{\star}\|_{P^{-1}}^{2}
   -2\eta\sum_{i=1}^{d}p_i\frac{1}{p_i}\partial_i f(w^{(t)})\bigl(w^{(t)}_i-w^{\star}_i\bigr)\\
   &\;+\eta^{2}\sum_{i=1}^{d}p_i\frac{1}{p_i^{2}}\partial_i f(w^{(t)})^{2}\\
=&\; \|w^{(t)}-w^{\star}\|_{P^{-1}}^{2}
   -2\eta\ip{\nabla f(w^{(t)})}{w^{(t)}-w^{\star}}\\
   &\;+\eta^{2}\sum_{i=1}^{d}\frac{1}{p_i}\partial_i f(w^{(t)})^{2}.
\end{aligned}
\tag{10}
\]

\noindent\textbf{[Step 5]}  Relate the last term to the smoothness constants.
For each coordinate, by Lemma \ref{lem:smooth} with $\alpha=-\eta\partial_i f(w^{(t)})/p_i$,
\[
f\!\bigl(w^{(t)}-\eta\frac{1}{p_i}\partial_i f(w^{(t)})e_i\bigr)
\le f(w^{(t)})-\eta\frac{1}{p_i}\partial_i f(w^{(t)})^{2}
   +\frac{L_i}{2}\eta^{2}\frac{1}{p_i^{2}}\partial_i f(w^{(t)})^{2}.
\]
Rearranging gives
\[
\frac{1}{p_i}\partial_i f(w^{(t)})^{2}
\le \frac{2}{\eta}\bigl[f(w^{(t)})-f(w^{(t)}-\eta\frac{1}{p_i}\partial_i f(w^{(t)})e_i)\bigr]
   +L_i\eta\frac{1}{p_i^{2}}\partial_i f(w^{(t)})^{2}.
\tag{11}
\]

\noindent\textbf{[Step 6]}  Sum \eqref{11} over $i$ and use that the update in the
algorithm changes only the sampled coordinate. Hence
\[
\sum_{i=1}^{d}\frac{1}{p_i}\partial_i f(w^{(t)})^{2}
\le \frac{2}{\eta}\bigl[f(w^{(t)})- \E\bigl[f(w^{(t+1)})\mid\mathcal{F}_t\bigr]\bigr]
   +\eta\sum_{i=1}^{d}L_i\frac{1}{p_i^{2}}\partial_i f(w^{(t)})^{2}.
\tag{12}
\]

\noindent\textbf{[Step 7]}  Plug \eqref{12} into \eqref{10}:
\[
\begin{aligned}
\E\!\bigl[\|w^{(t+1)}-w^{\star}\|_{P^{-1}}^{2}\mid\mathcal{F}_t\bigr]
\le &\; \|w^{(t)}-w^{\star}\|_{P^{-1}}^{2}
      -2\eta\ip{\nabla f(w^{(t)})}{w^{(t)}-w^{\star}} \\
    &\; +2\bigl[f(w^{(t)})-\E[f(w^{(t+1)})\mid\mathcal{F}_t]\bigr]
      +\eta^{2}\sum_{i=1}^{d}L_i\frac{1}{p_i^{2}}\partial_i f(w^{(t)})^{2}.
\end{aligned}
\tag{13}
\]

\noindent\textbf{[Step 8]}  Apply convexity (Assumption \ref{as:convex}) to bound
the inner product:
\[
f(w^{(t)})-f(w^{\star})\le \ip{\nabla f(w^{(t)})}{w^{(t)}-w^{\star}}.
\tag{14}
\]
Thus $-2\eta\ip{\nabla f(w^{(t)})}{w^{(t)}-w^{\star}}
\le -2\eta\bigl[f(w^{(t)})-f(w^{\star})\bigr]$.

\noindent\textbf{[Step 9]}  Insert \eqref{14} into \eqref{13} and rearrange:
\[
\begin{aligned}
2\eta\bigl[f(w^{(t)})-f(w^{\star})\bigr]
\le &\; \|w^{(t)}-w^{\star}\|_{P^{-1}}^{2}
      -\E\!\bigl[\|w^{(t+1)}-w^{\star}\|_{P^{-1}}^{2}\mid\mathcal{F}_t\bigr] \\
    &\; +2\bigl[f(w^{(t)})-\E[f(w^{(t+1)})\mid\mathcal{F}_t]\bigr]
      +\eta^{2}\sum_{i=1}^{d}L_i\frac{1}{p_i^{2}}\partial_i f(w^{(t)})^{2}.
\end{aligned}
\tag{15}
\]

\noindent\textbf{[Step 10]}  Drop the non‑negative term $2\bigl[f(w^{(t)})-\E f(w^{(t+1)})\mid\mathcal{F}_t\bigr]$
to obtain a simpler inequality:
\[
2\eta\bigl[f(w^{(t)})-f(w^{\star})\bigr]
\le \|w^{(t)}-w^{\star}\|_{P^{-1}}^{2}
    -\E\!\bigl[\|w^{(t+1)}-w^{\star}\|_{P^{-1}}^{2}\mid\mathcal{F}_t\bigr]
    +\eta^{2}\sum_{i=1}^{d}L_i\frac{1}{p_i^{2}}\partial_i f(w^{(t)})^{2}.
\tag{16}
\]

\noindent\textbf{[Step 11]}  Use the stepsize bound $\eta\le 1/L_{\max}$ and the fact that
$L_i\le L_{\max}$ for all $i$:
\[
\eta^{2}L_i\frac{1}{p_i^{2}}\partial_i f(w^{(t)})^{2}
\le \eta\frac{1}{p_i}\partial_i f(w^{(t)})^{2}.
\]
Summing over $i$ gives
\[
\eta^{2}\sum_{i=1}^{d}L_i\frac{1}{p_i^{2}}\partial_i f(w^{(t)})^{2}
\le \eta\sum_{i=1}^{d}\frac{1}{p_i}\partial_i f(w^{(t)})^{2}.
\tag{17}
\]

\noindent\textbf{[Step 12]}  Combine \eqref{16} and \eqref{17}:
\[
2\eta\bigl[f(w^{(t)})-f(w^{\star})\bigr]
\le \|w^{(t)}-w^{\star}\|_{P^{-1}}^{2}
    -\E\!\bigl[\|w^{(t+1)}-w^{\star}\|_{P^{-1}}^{2}\mid\mathcal{F}_t\bigr]
    +\eta\sum_{i=1}^{d}\frac{1}{p_i}\partial_i f(w^{(t)})^{2}.
\tag{18}
\]

\noindent\textbf{[Step 13]}  Observe that the last term on the right‑hand side of
\eqref{18} is exactly the expression that appeared in \eqref{10}. Re‑substituting
the equality from \eqref{10} (without the variance term) yields
\[
\eta\sum_{i=1}^{d}\frac{1}{p_i}\partial_i f(w^{(t)})^{2}
\le \|w^{(t)}-w^{\star}\|_{P^{-1}}^{2}
    -\E\!\bigl[\|w^{(t+1)}-w^{\star}\|_{P^{-1}}^{2}\mid\mathcal{F}_t\bigr]
    +2\eta\bigl[f(w^{(t)})-f(w^{\star})\bigr].
\]
Cancelling identical terms from both sides leaves
\[
0\le \|w^{(t)}-w^{\star}\|_{P^{-1}}^{2}
    -\E\!\bigl[\|w^{(t+1)}-w^{\star}\|_{P^{-1}}^{2}\mid\mathcal{F}_t\bigr].
\tag{19}
\]

\noindent\textbf{[Step 14]}  Take full expectation and telescope over $t=0,\dots,T-1$:
\[
\E\!\bigl[\|w^{(T)}-w^{\star}\|_{P^{-1}}^{2}\bigr]
\le \|w^{(0)}-w^{\star}\|_{P^{-1}}^{2}
    -2\eta\sum_{t=0}^{T-1}\E\!\bigl[f(w^{(t)})-f(w^{\star})\bigr].
\tag{20}
\]

\noindent\textbf{[Step 15]}  Since the norm is non‑negative, discard it from the left‑hand side and solve for the average suboptimality:
\[
\frac{1}{T}\sum_{t=0}^{T-1}\E\!\bigl[f(w^{(t)})-f(w^{\star})\bigr]
\le \frac{\|w^{(0)}-w^{\star}\|_{P^{-1}}^{2}}{2\eta T}.
\tag{21}
\]
Because $f$ is convex, the function value at the averaged iterate
$\bar w_T:=\frac1T\sum_{t=0}^{T-1}w^{(t)}$ does not exceed the average of the
function values, i.e.
\[
\E\!\bigl[f(\bar w_T)-f(w^{\star})\bigr]\le \frac{1}{T}\sum_{t=0}^{T-1}\E\!\bigl[f(w^{(t)})-f(w^{\star})\bigr].
\]
Combining with \eqref{21} yields the bound \eqref{5} of Theorem \ref{thm:conv}.
\hfill $\square$

\bigskip
\noindent\textbf{Linear rate under strong convexity.}  
When Assumption \ref{as:strong} holds, from convexity we have
\[
f(w^{(t)})-f(w^{\star})\ge \frac{\mu}{2}\|w^{(t)}-w^{\star}\|^{2}.
\]
Repeating the steps above but replacing the bound \eqref{14} with the strong
convexity inequality gives
\[
\|w^{(t+1)}-w^{\star}\|_{P^{-1}}^{2}
\le (1-2\eta\mu)\|w^{(t)}-w^{\star}\|_{P^{-1}}^{2},
\]
and unrolling yields
\[
\E\!\bigl[f(w^{(T)})-f(w^{\star})\bigr]
\le (1-\eta\mu)^{T}\bigl[f(w^{(0)})-f(w^{\star})\bigr],
\]
which is exactly \eqref{6}. \hfill $\blacksquare$

\section*{Rate Derivation (With Constants)}
For the convex case the bound \eqref{5) can be rewritten as
\[
\E\!\bigl[f(\bar w_T)-f(w^{\star})\bigr]
\le \frac{1}{2\eta T}\sum_{i=1}^{d}\frac{(w^{(0)}_i-w^{\star}_i)^{2}}{p_i}.
\]
If we choose $p_i\propto L_i^{-1/2}$, the denominator scales as $\sum_i\sqrt{L_i}$,
yielding the classic $O\!\bigl(\frac{\sum_i\sqrt{L_i}}{T}\bigr)$ rate.

\section*{Special Cases}
\begin{itemize}[leftmargin=*]
\item \textbf{Uniform probabilities.} $p_i=1/d$ gives $\|\,\cdot\,\|_{P^{-1}}^{2}=d\|\cdot\|^{2}$,
and the bound reduces to $O\!\bigl(d/(L_{\max}T)\bigr)$, matching standard RCD.
\item \textbf{Importance sampling.} Setting $p_i = \frac{L_i^{1/2}}{\sum_j L_j^{1/2}}$
minimises the constant $\sum_i \frac{1}{p_i}$ and improves practical performance.
\item \textbf{Block coordinates.} The same analysis extends by treating each
block as a single coordinate and using block‑wise smoothness constants.
\end{itemize}

\end{document}